\section{The cohomological Fourier--Mukai transform}\label{sec:fm}

With the character obstruction proved, we turn to the divisor obstruction of
\Cref{thm:divisor}. Its one nontrivial ingredient is a computation: the
cohomological Fourier--Mukai transform sends each power of the polarisation
class to a rational multiple of a power of the polarisation class on the
other side. We prove this in closed form. The computation takes place in a
cohomology ring which, together with the classes involved, depends up to
isomorphism only on the dimension, so one calculation covers every
principally polarised abelian variety at once. Since the Todd class of an
abelian variety is trivial, the passage from sheaves to cohomology is a
formal step.

\subsection{Grothendieck--Riemann--Roch on an abelian variety}

\begin{proposition}\label{prop:grr}
Let $A$ be an abelian variety of dimension $g$ and $\cG$ an object of
$D^{b}(\Av)$. Then
\[
  \ch\bigl(\Phi_{\cP}(\cG)\bigr)=\Phi_{\cP}\bigl(\ch(\cG)\bigr),
\]
where the transform on the right is the cohomological one,
$\Phi_{\cP}(x)=p_{*}\bigl(q^{*}x\cdot\ch(\cP)\bigr)$.
\end{proposition}

\begin{proof}
Grothendieck--Riemann--Roch for the projection
$p\colon A\times\Av\to A$, a proper morphism of smooth varieties, reads
\[
  \ch\bigl(Rp_{*}\cK\bigr)\cdot\td(T_{A})
  =p_{*}\bigl(\ch(\cK)\cdot\td(T_{A\times\Av})\bigr)
\]
for $\cK\in D^{b}(A\times\Av)$. The tangent bundle of an abelian variety is
trivial, so $\td(T_{A})=1$ and $\td(T_{A\times\Av})=1$. Take
$\cK=Lq^{*}\cG\otimes^{L}\cP$, so that $Rp_{*}\cK=\Phi_{\cP}(\cG)$; the
multiplicativity of the Chern character and $\ch(Lq^{*}\cG)=q^{*}\ch(\cG)$
give the statement.
\end{proof}

\subsection{The universal model}

\begin{notation}\label{not:model}
Let $A$ be a principally polarised abelian variety of dimension $g$ with
polarisation class $\theta$, and put $V=H^{1}(A,\QQ)$, a $\QQ$-vector space of
dimension $2g$ carrying the symplectic form attached to $\theta$. Choose a
symplectic basis $e_{1},\dots,e_{2g}$ of $V$ for this form, so that
\begin{equation}\label{eq:thetamodel}
  \theta=\sum_{j=1}^{g}e_{j}\wedge e_{g+j}.
\end{equation}
The principal polarisation identifies $\Av$ with $A$ and
$H^{1}(\Av,\QQ)$ with the dual $V^{\vee}$; let $f_{1},\dots,f_{2g}$ be the
basis dual to $e_{1},\dots,e_{2g}$, so that
\begin{equation}\label{eq:thetapmodel}
  \theta'=\sum_{j=1}^{g}f_{j}\wedge f_{g+j}.
\end{equation}
Then $H^{*}(A\times\Av,\QQ)=\bigwedge^{*}(V\oplus V^{\vee})$.
\end{notation}

\begin{lemma}[The Poincar\'e class]\label{lem:poincare}
In the model of \Cref{not:model},
\begin{equation}\label{eq:ell}
  \ell\coloneqq c_{1}(\cP)=\sum_{i=1}^{2g}e_{i}\wedge f_{i}
  \;\in\;V\otimes V^{\vee}\subset\textstyle\bigwedge^{2}(V\oplus V^{\vee}),
\end{equation}
the canonical element of $V\otimes V^{\vee}=\Hom(V,V)$ corresponding to the
identity, and $\ch(\cP)=e^{\ell}$.
\end{lemma}

\begin{proof}
The Poincar\'e bundle is normalised to be trivial on $A\times\{0\}$ and on
$\{0\}\times\Av$. The components of $c_{1}(\cP)$ in $\bigwedge^{2}V$ and in
$\bigwedge^{2}V^{\vee}$ are its restrictions to these two subvarieties, so
they vanish, and $c_{1}(\cP)$ lies in the middle summand $V\otimes V^{\vee}$
of
$\bigwedge^{2}(V\oplus V^{\vee})=\bigwedge^{2}V\oplus(V\otimes V^{\vee})\oplus\bigwedge^{2}V^{\vee}$.
The defining property of $\cP$ is that its restriction to $A\times\{\lambda\}$
is the line bundle $\lambda$; on cohomology this says that the induced map
$H^{1}(\Av,\QQ)\to H^{1}(A,\QQ)^{\vee}$ obtained by contracting $c_{1}(\cP)$ is
the canonical identification. Equivalently, the element of $\Hom(V,V)$
determined by $c_{1}(\cP)$ is the identity, which is \eqref{eq:ell}. Since
$\cP$ is a line bundle, $\ch(\cP)=e^{c_{1}(\cP)}$.
\end{proof}

\begin{remark}\label{rem:universal}
Every principally polarised abelian variety of dimension $g$ gives, up to
isomorphism, the same data in \Cref{not:model}: a symplectic space of
dimension $2g$ over $\QQ$, its dual, the symplectic form, and the canonical
element. The constants computed
below therefore depend only on $g$ and $k$, and not on $A$.
\end{remark}

\subsection{The closed form}

\begin{theorem}[Transform of the polarisation powers]\label{thm:fmtheta}
Let $A$ be a principally polarised abelian variety of dimension $g$. Then for
every $0\le k\le g$,
\begin{equation}\label{eq:fmtheta}
  \Phi_{\cP}\bigl(\theta'^{\,k}\bigr)
  \;=\;(-1)^{\,g(g+1)/2\,+\,k}\,\frac{k!}{(g-k)!}\;\theta^{\,g-k},
\end{equation}
and $\Phi_{\cP}(\theta'^{\,k})=0$ for $k>g$. In particular
$\Phi_{\cP}\bigl(\QQ[\theta']\bigr)=\QQ[\theta]$.
\end{theorem}

\begin{proof}
Work in the model of \Cref{not:model}. By definition
$\Phi_{\cP}(x)=p_{*}(q^{*}x\cdot e^{\ell})$, and $p_{*}$ extracts the
coefficient of the top exterior form $f_{1}\wedge\cdots\wedge f_{2g}$ in the
$V^{\vee}$ variables: it sends $y\wedge f_{1}\wedge\cdots\wedge f_{2g}$ to
$y$ for $y\in\bigwedge^{*}V$, and kills every class of lower degree in the
$V^{\vee}$ variables.

The classes $e_{i}\wedge f_{i}$ have even degree, so they commute, and they
have square zero. Hence for every $m$,
\begin{equation}\label{eq:ellpower}
  \frac{\ell^{m}}{m!}
  =\sum_{\substack{S\subseteq\{1,\dots,2g\}\\ |S|=m}}\ \prod_{i\in S}e_{i}\wedge f_{i},
\end{equation}
the product being taken in increasing order of $i$. Likewise, from
\eqref{eq:thetapmodel},
\begin{equation}\label{eq:thetappower}
  \frac{\theta'^{\,k}}{k!}
  =\sum_{\substack{T\subseteq\{1,\dots,g\}\\ |T|=k}}\ \prod_{j\in T}f_{j}\wedge f_{g+j},
  \qquad
  \frac{\theta^{\,g-k}}{(g-k)!}
  =\sum_{\substack{U\subseteq\{1,\dots,g\}\\ |U|=g-k}}\ \prod_{j\in U}e_{j}\wedge e_{g+j}.
\end{equation}

Fix $k$. In the product $\theta'^{\,k}\cdot e^{\ell}$, the terms that survive
$p_{*}$ are those of total degree $2g$ in the $f$ variables. A term of
\eqref{eq:thetappower} indexed by $T$ carries the $2k$ variables
$\{f_{j},f_{g+j}:j\in T\}$, and a term of \eqref{eq:ellpower} indexed by $S$
carries the $|S|$ variables $\{f_{i}:i\in S\}$. For the product to be nonzero
the two index sets must be disjoint, and for the total $f$-degree to be $2g$
we need
\[
  S=\{1,\dots,2g\}\setminus\bigl(T\cup(g+T)\bigr),
  \qquad |S|=2g-2k .
\]
So $S$ is determined by $T$, only the term $\ell^{2g-2k}/(2g-2k)!$ of
$e^{\ell}$ contributes, and the surviving terms are in bijection with the
$k$-subsets $T$ of $\{1,\dots,g\}$. Writing $U=\{1,\dots,g\}\setminus T$ we
have $S=U\cup(g+U)$, so the $V$-part of the term indexed by $T$ is
$\prod_{i\in U\cup(g+U)}e_{i}$, which up to sign is the term of
\eqref{eq:thetappower} indexed by $U$. The bijection $T\leftrightarrow U$ is
the one between the terms of $\theta'^{\,k}/k!$ and those of
$\theta^{\,g-k}/(g-k)!$.

It remains to compute the sign $\sigma$ with which each term contributes, and
to check that it is the same for all $T$. The sign is that of the
rearrangement of the ordered word
\[
  \Bigl(\prod_{j\in T}f_{j}f_{g+j}\Bigr)\cdot\Bigl(\prod_{i\in S}e_{i}f_{i}\Bigr)
\]
into the target word
$\bigl(\prod_{j\in U}e_{j}e_{g+j}\bigr)\cdot\bigl(f_{1}\cdots f_{2g}\bigr)$,
which uses the same generators. Relabelling the $g$ symplectic pairs by a
permutation $\pi$ of $\{1,\dots,g\}$ is an automorphism of the exterior
algebra. It
carries the word for $T$ to the word for $\pi(T)$, because the blocks
$f_{j}f_{g+j}$ and $e_{i}f_{i}$ have degree two and may be reordered freely,
and it carries the target word for $U$ to the target word for $\pi(U)$,
because reordering
$f_{\pi(1)}\cdots f_{\pi(g)}f_{g+\pi(1)}\cdots f_{g+\pi(g)}$ costs the
square of the sign of $\pi$. Since the permutations of $\{1,\dots,g\}$ act
transitively on the $k$-subsets, $\sigma=\sigma(g,k)$, and we evaluate it at
$T=\{1,\dots,k\}$. Put $r=g-k$. Moving the $2r$ variables $e_{i}$ of the
second factor to its front costs $(-1)^{r(2r-1)}=(-1)^{r}$, moving them
further past the $2k$ variables of the first factor costs nothing, and
arranging them as $\prod_{j\in U}e_{j}e_{g+j}$ costs $(-1)^{r(r-1)/2}$. The
$f$ variables are left in the order
$f_{1}f_{g+1}\cdots f_{k}f_{g+k}\,f_{k+1}\cdots f_{g}\,f_{g+k+1}\cdots f_{2g}$,
and bringing them to $f_{1}\cdots f_{2g}$ costs $(-1)^{k(k-1)/2+rk}$. As
$r+r(r-1)/2+k(k-1)/2+rk=g(g+1)/2-k$, this gives
\begin{equation}\label{eq:sign}
  \sigma(g,k)=(-1)^{\,g(g+1)/2\,+\,k}.
\end{equation}
Assembling, and accounting for the factorials in \eqref{eq:thetappower},
\[
  \Phi_{\cP}(\theta'^{\,k})
  =k!\cdot\sigma(g,k)\cdot\frac{\theta^{\,g-k}}{(g-k)!},
\]
which is \eqref{eq:fmtheta}. For $k>g$ we have $\theta'^{\,k}=0$, because
$\theta'^{\,k}$ lies in $\bigwedge^{2k}V^{\vee}$, which vanishes for
$2k>2g$. The last assertion follows because the constants in
\eqref{eq:fmtheta} are nonzero for every $0\le k\le g$.
\end{proof}

\begin{remark}[The sign]\label{rem:signcheck}
The sign \eqref{eq:sign} depends only on $g$ and $k$, by the transitivity
argument in the proof, and the count at $T=\{1,\dots,k\}$ fixes it. For even $g=2n$, the case of Weil type,
$g(g+1)/2=n(2n+1)$ has the parity of $n$, and \eqref{eq:sign} reads
$\sigma(2n,k)=(-1)^{n+k}$.
\end{remark}

\begin{corollary}[Stability of the polarisation subring]\label{cor:stability}
Let $A$ be a principally polarised abelian variety and let
$x\in H^{*}(\Av,\QQ)$ lie in the subring $\QQ[\theta']$ generated by $\theta'$.
Then $\Phi_{\cP}(x)\in\QQ[\theta]$.
\end{corollary}

\begin{proof}
The ring $\QQ[\theta']$ is spanned over $\QQ$ by the powers
$\theta'^{\,k}$, \Cref{thm:fmtheta} sends each of them to a rational
multiple of a power of $\theta$, and $\Phi_{\cP}$ is linear.
\end{proof}

\begin{remark}[Non-principal polarisations]\label{rem:nonprincipal}
If $\theta$ is a polarisation of degree $D^{2}$ rather than a principal
one, the isogeny $\phi_{\theta}\colon A\to\Av$ has degree $D^{2}$,
and $\phi_{\theta}^{*}\theta'=D\theta$ for a suitable normalisation of
$\theta'$. The argument of \Cref{thm:fmtheta} goes through with $\theta$
replaced by $D^{-1}\phi_{\theta}^{*}\theta'$, and the constants change
by powers of $D$, all of which are nonzero rationals. The statements below
concern only membership in $\QQ[\theta]$, not the values of the constants, so
we state them for principal polarisations and use them in general.
\end{remark}
